\documentclass[ ../main.tex]{subfiles}
\providecommand{\mainx}{..}

\begin{document}
\section{Rate-distorted $(1-\alpha)$-hashing}







We generalize the perfect hash function by relaxing the require on no collisions;
instead, we allow collisions and place constraints on the maximum space size of the data that represents the perfect hash function.
The result is a $(1-\alpha)$-hash function where proportion $\alpha$ of the elements in the objective set failed to be perfectly hashed.
The \emph{perfect hash function} is a special case where $\alpha=0$.

The number of elements that collide on any particular trial is binomially distributed,
\begin{equation}
	\RV{T} \sim B(m,p)
\end{equation}
where $p = \fprate$.

There are a few different ways of treating $\alpha$.
One way is to treat $\alpha$ as a random variable that is induced by a maximum bit rate per element, e.g., on average the maximum number of bits per element is $b$.
Thus, if there are $m$ elements, the constructor is limited to finding a bit string of length $b m$.

Another way to proceed is to say that we are satisfied with proportion $(1-\alpha)$ elements or more being perfectly hashed.


The probability that 
\begin{equation}
	\Prob{\RV{Q}_t \leq q \Given \mathlarger{\cap}_{u=t+1}^{m} (\RV{Q}_{u} > q)}
\end{equation}



The \emph{data structure} for the rate-distorted perfect hash function with a maximum bit length $b$ is defined as $\PH(b) = \BitSet^{b} \times \NatSet \times \RatSet$.
An optimal constructor for $\PH(b)$ with a load factor $r$,
\begin{equation}
	\ph \colon \PS{\Set{U}} \times \RatSet \times \NatSet \mapsto \PH(b)\,,
\end{equation}
is given by
\begin{equation}
	\ph(\Set{X}, r, b) \coloneqq \Tuple{n',N, \alpha}
\end{equation}
where
\begin{equation}
\begin{split}
\ell		& \coloneqq 2^b\,, m \coloneqq \Card{\Set{X}}\,, N \coloneqq m / r\,, k \coloneqq \lceil \log_2 N \rceil\,,\\
\beta(x,n) 	& \coloneqq \Fun{pad}\left(\hash(x' \cat n') \mod N, k\right)\,,\\
\Set{Y}[n] 	& \coloneqq \SetBuilder{\beta(x,n) \in \BitSet^k}{x \in \Set{X}}\,,\\
\gamma(k)	& \coloneqq \SetBuilder{ j \in \NatSet_{\leq \ell} }{ \Set{Y}[j] \in \PS{\BitSet^N} \land \Card{\Set{Y}[j]} = k}\,,\\
p			& \coloneqq \argmax_j \SetBuilder{ j }{\gamma(j) \neq \EmptySet}\,\\
\alpha		& \coloneqq 1-\frac{p}{m}\,\\
n			& \coloneqq \min \gamma(p)\,.
\end{split}
\end{equation}




\end{document}




%\begin{proof}
%	Coupon collector problem.
%	$\Card{\Set{S}} = m$
%	$N = \frac{m}{r}$
%	\begin{equation}
%		\RV{K_j} \sim \geodist\!\left(\frac{N-1}{N}\right)
%	\end{equation}
%	$\Pr[\RV{K_j} \leq k] = 1 - \left(1 - \frac{N-j+1}{N}\right)^{k+1}$
%	A maximum of $k$ hashes...
%	$p(k) = \Pr[\RV{K_1} \leq k \cap \cdots \cap \RV{K_m} \leq k] = \prod_{j=1}^{m}\left(1 - \left(\frac{j-1}{N}\right)^{k+1}\right)$
%	$p(k) = \Pr[\RV{K_1} \leq k \cap \cdots \cap \RV{K_m} \leq k] = \prod_{j=1}^{m-1}\left(1 - \left(\frac{j}{N}\right)^{k+1}\right)$
%	$p(k) = \Pr[\RV{K_1} \leq k \cap \cdots \cap \RV{K_m} \leq k] = \prod_{j=1}^{m-1}\frac{N^{k+1} - j^{k+1}}{N^{k+1}}$
%	$p(k) = \Pr[\RV{K_1} \leq k \cap \cdots \cap \RV{K_m} \leq k] = N^{-(m-1)(k+1)} \prod_{j=1}^{m-1}\left(N^{k+1} - j^{k+1}\right)$
%\end{proof}


